\subsection{Code chi tiết cho Section 6}\label{phu-luc-code-section-6}

\noindent\textbf{Chi tiết quy trình kiểm toán thủ công cho Section 6}

\begin{enumerate}
    \item \textbf{6.1 -- Tránh tích tụ quá nhiều Docker Images dư thừa (Image sprawl):}
    Đếm và đối chứng số lượng image:
{\footnotesize
\begin{verbatim}
# Số lượng images đang thực sự được sử dụng bởi các container
docker ps -qa | xargs -r docker inspect --format "{{.Image}}" \
  | sort -u | wc -l

# Tổng số images hiện có trên máy chủ
docker image list -a -q | wc -l

# Liệt kê chi tiết để rà soát
docker image list --format \
  'table {{.Repository}}:{{.Tag}}\t{{.ID}}\t{{.Size}}'
\end{verbatim}
}

    \item \textbf{6.2 -- Tránh tích tụ quá nhiều Docker Containers đã dừng (Container sprawl):}
    Rà soát các container dư thừa:
{\scriptsize
\begin{verbatim}
# Đọc thống kê containers từ hệ thống
docker info --format \
  'Containers: {{.Containers}},
   Running: {{.ContainersRunning}},
   Stopped: {{.ContainersStopped}}'

# Liệt kê tất cả container kèm theo trạng thái hoạt động của chúng
docker ps -a --format \
  'table {{.ID}}\t{{.Image}}\t{{.Status}}\t{{.Names}}'
\end{verbatim}
}
\end{enumerate}

\noindent\textbf{Chi tiết quy trình khắc phục thủ công cho Section 6}

\begin{enumerate}
    \item \textbf{6.1 -- Loại bỏ Docker images dư thừa (Image sprawl):}
    Dọn dẹp các images không còn được sử dụng bởi bất kỳ container nào:
{\footnotesize
\begin{verbatim}
# Phương án 1: Dọn dẹp triệt để toàn bộ image không sử dụng
docker image prune -a -f

# Phương án 2: Dọn dẹp có chọn lọc (giữ lại các image mới hơn 30 ngày)
docker image prune -a --filter "until=720h" -f

# Thiết lập cron job để tự động dọn dẹp hàng tuần vào 3h sáng Chủ Nhật
echo '0 3 * * 0 root docker image prune -a --filter "until=720h" -f >> /var/log/docker-prune.log 2>&1' | \
  sudo tee -a /etc/cron.d/docker-cleanup
\end{verbatim}
}

    \item \textbf{6.2 -- Loại bỏ các Docker container đã dừng (Container sprawl):}
    Dọn dẹp các container mồ côi đã kết thúc phiên chạy:
{\footnotesize
\begin{verbatim}
# Xóa bỏ tất cả các container đang ở trạng thái dừng
docker container prune -f

# Phòng ngừa: Khuyến khích sử dụng cờ --rm cho container chạy kiểm thử (one-off)
docker run --rm <image_name> <command>

# Thiết lập cron job tự động dọn container dừng cũ hơn 48 giờ hàng ngày lúc 4h sáng
echo '0 4 * * * root docker container prune --filter "until=48h" -f >> /var/log/docker-prune.log 2>&1' | \
  sudo tee -a /etc/cron.d/docker-cleanup
\end{verbatim}
}
\end{enumerate}

\noindent\textbf{Chi tiết kịch bản kiểm toán tự động cho Section 6}

Kịch bản kiểm toán tự động chống image sprawl và container sprawl:

{\footnotesize
\begin{verbatim}
# Kiểm toán 6.1: Tránh tích tụ quá nhiều image (image sprawl)
- name: "6.1 | Count images actively used by running containers"
  ansible.builtin.shell:
    cmd: |
      docker ps -qa | xargs -r docker inspect --format "{{.Image}}" | sort -u | wc -l
  register: active_images_command
  changed_when: false
  failed_when: false

- name: "6.1 | Count all images on host"
  ansible.builtin.shell:
    cmd: "docker image list -a -q | wc -l"
  register: all_images_command
  changed_when: false
  failed_when: false

# Kiểm toán 6.2: Tránh tích tụ quá nhiều container dừng (container sprawl)
- name: "6.2 | Get total container count"
  ansible.builtin.shell:
    cmd: "docker info --format '{{ .Containers }}'"
  register: total_container_command
  changed_when: false
  failed_when: false

- name: "6.2 | Get running container count"
  ansible.builtin.shell:
    cmd: "docker info --format '{{ .ContainersRunning }}'"
  register: total_running_container_command
  changed_when: false
  failed_when: false

- name: "6.2 | Get stopped container count"
  ansible.builtin.shell:
    cmd: "docker info --format '{{ .ContainersStopped }}'"
  register: total_not_running_container_command
  changed_when: false
  failed_when: false
\end{verbatim}
}
